\subsection{Lifting of paths}

PROPOSITION 3. --- Let B be a topological space and let \(p\colon E \to B\) be a covering. Denote by \(\widetilde{p}\colon \mathcal{C}_{c}(\mathbf{I};E) \to \mathcal{C}_{c}(\mathbf{I};B)\) the map \(c \mapsto p \circ c\) . Denote by \(o_{\mathrm{E}}\colon \mathcal{C}_{c}(\mathbf{I};E) \to E\) (resp. \(o_{\mathrm{B}}\colon \mathcal{C}_{c}(\mathbf{I};B) \to B\) ) the map defined by \(c \mapsto c(0)\) . Then the diagram

\[
\begin{array}{c} \mathcal {C} _ {c} (\mathbf {I}; \mathrm{E}) \xrightarrow {o _ {\mathrm{E}}} \mathrm{E} \\ \Big \downarrow \widetilde {p} \\ \mathcal {C} _ {c} (\mathbf {I}; \mathrm{B}) \xrightarrow {o _ {\mathrm{B}}} \mathrm{B} \end{array}
\]

is a Cartesian square.

The topological space I is locally connected, locally compact, and simply connected. The proposition thus follows from prop. 9 of I, p.~131, applied with T = I and t = 0.

COROLLARY 1. --- The map \(\widetilde{p}\colon\mathcal{C}_{c}(\mathbf{I};\mathrm{E})\to\mathcal{C}_{c}(\mathbf{I};\mathrm{B})\) is a covering.

This follows from I, p.~71, cor. 2 of prop. 1.

COROLLARY 2 (Lifting of paths). --- Let \(p\) : E → B be a covering, let x be a point of E and a = p(x). The map c ↦ p ∘ c is a homeomorphism from the space \(\Lambda_x\) (E) of paths in E with origin x onto the space \(\Lambda_a\) (B) of paths in B with origin a.

With the notations of proposition 3, we have \(\Lambda_{a}(\mathrm{B}) = o_{\mathrm{B}}^{-1}(a)\) and \(\Lambda_x(\mathrm{E}) = o_{\mathrm{E}}^{-1}(x)\) . The corollary then follows from I, p.~10, cor. of prop. 4.

We shall say that a continuous map \(p \colon \mathrm{E} \to \mathrm{B}\) satisfies the path-lifting property if, for every path \(c \colon \mathbf{I} \to \mathbf{B}\) and every point \(x\) of E lifting \(c(0)\) , there exists a path \(c'\) in E with origin \(x\) which lifts \(c\) . Such a path is then unique if \(p\) is étale and separated (I, p.~34, cor. 1 of prop. 11). Let E be a covering and \(p\) its projection; the map \(p\) is étale, separated, and possesses the path-lifting property (corollary 2). For a partial converse, cf.~III, p.~315, corollary.

PROPOSITION 4. --- Let \(p \colon \mathrm{E} \to \mathrm{B}\) be an étale and separated map which satisfies the path-lifting property. Let \(a\) and \(b\) be points of \(\mathrm{B}\) and let \(x\) be a point of \(\mathrm{E}\) such that \(p(x) = a\) . Let \(c_0\) and \(c_1\) be two strictly homotopic paths joining \(a\) to \(b\) in \(\mathrm{B}\) . The paths with origin \(x\) in \(\mathrm{E}\) which respectively lift \(c_0\) and \(c_1\) have the same endpoint and are strictly homotopic.

Let \(\sigma\colon\mathbf{I}\times\mathbf{I}\to\mathbf{B}\) be a strict homotopy joining \(c_{0}\) to \(c_{1}\) . For \(s\in\mathbf{I}\) , let \(c_{s}^{\prime}\) be the unique path with origin \(x\) in E which lifts the path \(t\mapsto\sigma(t,s)\) . For \((t,s)\in\mathbf{I}\times\mathbf{I}\) , set \(\sigma^{\prime}(t,s)=c_{s}^{\prime}(t)\) ; we have \(p\circ\sigma^{\prime}=\sigma\) . The map \(\sigma^{\prime}\) is constant on \(\{0\}\times\mathbf{I}\) ; by construction, it is continuous on \(\mathbf{I}\times\{s\}\) for all \(s\in\mathbf{I}\) . It is therefore continuous, by I, p.~37, cor. 1 of th. 1. In particular, the map from \(\mathbf{I}\) to \(\mathrm{E}_{b}\) given by \(s\mapsto\sigma^{\prime}(1,s)\) is a continuous lifting of the constant path with image \(b\) ; it is therefore constant. This proves that the map \(\sigma^{\prime}\) is a strict homotopy joining \(c_{0}^{\prime}\) to \(c_{1}^{\prime}\) .

COROLLARY 1. --- Let B be a topological space and let E be a covering of B; denote by p its projection. For every pair \((x,y)\) of points of E, the map \(\varpi_{x,y}(p)\colon\varpi_{x,y}(\mathrm{E})\to\varpi_{p(x),p(y)}(\mathrm{B})\) is injective. In particular, for every point x of E, the homomorphism \(\pi_{1}(p,x)\colon\pi_{1}(\mathrm{E},x)\to\pi_{1}(\mathrm{B},p(x))\) is injective.

COROLLARY 2. --- Let B be a topological space and let E be a covering of B; denote by p its projection. Let x be a point of E, set \(b = p(x)\) . For the class in \(\pi_{1}(B, b)\) of a loop c of B at a to belong to the image subgroup of the homomorphism \(\pi_{1}(p, x)\) , it is necessary and sufficient that the path c' with origin x lifting c be a loop of E at x.

The condition is obviously sufficient. Conversely, let \(c_{1}^{\prime}\) be a loop of E at x. Set \(c_{1} = p \circ c_{1}^{\prime}\) and suppose that the loops c and \(c_{1}\) are strictly homotopic. By prop. 4, the path \(c^{\prime}\) has the same endpoint as the path \(c_{1}^{\prime}\) . It is therefore a loop at x.

COROLLARY 3. --- Let B be a topological space and let \((\mathrm{E}_{1}, p_{1})\) and \((\mathrm{E}_{2}, p_{2})\) be coverings of B. Denote \((\mathrm{E}, p)\) their B-product space and let \(x = (x_{1}, x_{2})\) be a point of E. The image of the homomorphism \(\pi_{1}(p, x)\) is the intersection of the images of the homomorphisms \(\pi_{1}(p_{1}, x_{1})\) and \(\pi_{1}(p_{2}, x_{2})\) .

Let \(c\) be a loop in B based at \(p_{1}(x_{1}) = p_{2}(x_{2})\) and, for \(i \in \{1, 2\}\) , let \(c_{i}^{\prime}\) be the path with origin \(x_i\) in \(E_i\) lifting \(c\). The B-space \(E = E_{1} \times_{B} E_{2}\) is a covering of B (I, p.~72, cor. 3 of prop. 2), and the path \(c^{\prime}: t \mapsto (c_{1}^{\prime}(t), c_{2}^{\prime}(t))\) is the unique path with origin \(x\) in E lifting \(c\). For \(c^{\prime}\) to be a loop, it is necessary and sufficient that \(c_{1}^{\prime}\) and \(c_{2}^{\prime}\) be loops. Corollary 3 thus follows from Corollary 2.

